\documentclass[11pt,a4paper]{article}
\usepackage[T1]{fontenc}
\usepackage[utf8]{inputenc}
\usepackage{lmodern,amsmath,amssymb,amsthm,booktabs}
\usepackage[margin=27mm]{geometry}
\usepackage[hidelinks]{hyperref}
\hypersetup{pdftitle={Physical-Time Action Relaxation in an Elastic Collision Bath},
pdfauthor={navstokgap research working notes}}
\setlength{\emergencystretch}{3em}
\newtheorem{proposition}{Proposition}
\newcommand{\E}{\mathbb E}
\newcommand{\Var}{\operatorname{Var}}
\newcommand{\Cov}{\operatorname{Cov}}
\title{Physical-Time Action Relaxation in an Elastic Collision Bath}
\author{navstokgap research working notes}
\date{9 September 2026}
\begin{document}
\maketitle
\begin{abstract}
An elastic repeated-collision model produces an action-valued covariance
observable that relaxes exponentially in physical time to a positive constant.
We derive its rate and plateau directly from the masses, bath velocity variance
and collision frequency. Bounded bath velocities keep every collision below a
fixed speed ceiling. The exact parameter dependence isolates the additional
law needed for a universal action scale: energy and momentum conservation
determine the collision map, while the reservoir statistics set the plateau.
An equal-mass two-speed gas supplies a spatial realization: independent
Poisson spacings generate reversal rate $\rho u$ and plateau $mu/\rho$.
\end{abstract}

\section{A momentum receiver and a classical stochastic clock}
A tracer of mass $m\ge M>0$ meets fresh bath particles of mass $M$ on the
line. Collisions occur at an independent Poisson rate $\nu>0$. Incoming bath
velocities $W_n$ are independent, identically distributed, centered and obey
$|W_n|\le w<c/3$, with variance $s^2>0$. Initial tracer velocity is independent
of the bath and clock, centered, and bounded by $w$. Between collisions it is
constant; $X(t)=X(0)+\int_0^t V(r)\,dr$.

The nontrivial one-dimensional elastic collision map is
\begin{equation}\label{eq:map}
 V'=aV+bW,\quad W'=\frac{2m}{m+M}V-aW,\qquad
 a=\frac{m-M}{m+M},\quad b=\frac{2M}{m+M}=1-a.
\end{equation}
Direct substitution gives $mV'+MW'=mV+MW$ and
$m(V')^2+M(W')^2=mV^2+MW^2$.
Since $0\le a<1$, the tracer update is a convex combination and preserves
$|V|\le w$. The departing particle obeys
$|W'|\le(3m-M)w/(m+M)<3w<c$.

The model specifies a refreshed reservoir and a velocity-independent clock.
These are Barkai's collision-model premises
\cite[pp.~3--5, (2)--(12)]{Barkai2003StableEquilibrium}, specialized here to a
bounded bath and $m\ge M$.
Each collision accounts for momentum and energy transfer; the complete stream
supplies fresh particles and removes departing ones. A spatial gas realization
would additionally derive encounter rates and incoming velocity statistics.
The speed ceiling here is a kinematic restriction on Newtonian collisions in
the reservoir frame, rather than a Lorentz-covariant collision law.

\section{Stationary covariance and the action plateau}
Define the two inverse-time rates
\begin{equation}\label{eq:rates}
 \gamma=\nu(1-a)=\frac{2\nu M}{m+M},\qquad
 \beta=\nu(1-a^2)=\frac{4\nu mM}{(m+M)^2}.
\end{equation}
The velocity generator acts on bounded functions as
$Lf(v)=\nu\{\E[f(av+bW)]-f(v)\}$.

\begin{proposition}[Collision plateau]\label{prop:plateau}
The stationary velocity law is unique on $[-w,w]$ and is represented by the
convergent series $V_*\overset{d}=b\sum_{j=0}^\infty a^jW_j$.
Its variance and autocovariance are
\begin{equation}\label{eq:cov}
 S_*:=\Var V_*=\frac{b^2s^2}{1-a^2}=\frac{M}{m}s^2,
 \qquad C(r)=S_*e^{-\gamma r}.
\end{equation}
Thus the stationary displacement observable
$\mathsf h(\Delta)=m\Var[X(t+\Delta)-X(t)]/\Delta$ satisfies
\begin{equation}\label{eq:h}
 \mathsf h(\Delta)=H_*\left[1-\frac{1-e^{-\gamma\Delta}}
 {\gamma\Delta}\right],\qquad
 H_*:=\frac{2mS_*}{\gamma}=\frac{(m+M)s^2}{\nu}>0.
\end{equation}
\end{proposition}
\begin{proof}
The series converges absolutely and stays in $[-w,w]$, since
$b\sum a^j=1$ (for $a=0$ it is simply $W_0$).
Separating its first term proves invariance under one collision. Coupling two
initial velocities to the same clock and bath multiplies their difference by
$a^{N_t}$. Its mean contraction factor is
$\E a^{N_t}=e^{-\nu(1-a)t}$, also for $a=0$ with $0^0=1$.
Hence invariant probability laws on the bounded interval coincide.
Independence and centering give the variance in \eqref{eq:cov}.
Conditioning on the collision count yields
$\E[V(t+r)\mid V(t)]=e^{-\gamma r}V(t)$; stationary covariance follows.
Integrating the covariance over the displacement square gives
$\mathsf h(\Delta)=2m\int_0^\Delta(1-r/\Delta)C(r)\,dr$,
which evaluates to \eqref{eq:h}.
\end{proof}

This proof uses a contracting affine collision update; it requires neither a
finite velocity state space nor detailed balance of the refreshed bath.
The energy-valued reservoir parameter $\Theta:=Ms^2$ gives
$mS_*=\Theta$ and $H_*=\Theta(m+M)/(\nu M)$.
For a general bounded bath, $\Theta$ denotes twice the mean incoming kinetic
energy, without assuming a thermodynamic equilibrium distribution.
The stationary covariance and Green--Kubo integral also follow from
Barbier and Trizac's field-free Maxwell-kernel specialization
\cite[p.~19, (89)--(98)]{BarbierTrizac2012Field}. Their kernel exponent
called $\nu$ is zero there; our $\nu$ is a positive collision frequency.

\section{An observable relaxing in physical time}
The integrated future-lag covariance defines a second operational quantity:
\begin{equation}\label{eq:field}
 \mathsf a(t):=2m\int_0^\infty\Cov[V(t+r),V(t)]\,dr.
\end{equation}
It has action units and uses physical preparation time $t$, with the lag
integrated out. Its definition differs from the finite-duration displacement
observable; both share the same stationary long-duration plateau.

\begin{proposition}[Physical-time relaxation]\label{prop:relax}
For the centered initial ensemble specified above, with $S_0=\Var V(0)$,
\begin{align}
 S(t)&=S_*+(S_0-S_*)e^{-\beta t},\label{eq:variance}\\
 \mathsf a(t)&=\frac{2mS(t)}\gamma
 =H_*+(\mathsf a(0)-H_*)e^{-\beta t},\qquad
 \dot{\mathsf a}=\beta(H_*-\mathsf a).\label{eq:relax}
\end{align}
In particular a tracer prepared at rest has
$\mathsf a(t)=H_*(1-e^{-\beta t})>0$ for every $t>0$ and
$\mathsf a(t)\to H_*>0$ as $t\to\infty$.
\end{proposition}
\begin{proof}
The generator gives $Lv=-\gamma v$ and
$L(v^2)=-\beta v^2+\nu b^2s^2$. Centering is preserved, so
$\dot S=-\beta S+\nu b^2s^2=\beta(S_*-S)$.
The conditional mean from Proposition~\ref{prop:plateau} gives, even away from
stationarity, $\Cov[V(t+r),V(t)]=e^{-\gamma r}S(t)$.
Integration proves \eqref{eq:relax}.
\end{proof}

The original displacement definition retains a distinct duration parameter:
\begin{equation}\label{eq:window}
 \mathsf h_\Delta(t)=\frac{2m}{\gamma\Delta}
 \int_0^\Delta S(t+r)\bigl[1-e^{-\gamma(\Delta-r)}\bigr]\,dr.
\end{equation}
Indeed, for $r\le s$ the covariance inside its double integral is
$e^{-\gamma(s-r)}S(t+r)$; integrating over $s$ proves the formula.
At fixed $\Delta$, it approaches the stationary value in \eqref{eq:h}
exponentially at rate $\beta$. At every fixed $t$ it tends to zero as
$\Delta\downarrow0$, with $\mathsf h_\Delta(t)\sim mS(t)\Delta$ when
$S(t)>0$. Its long-duration limit is $H_*$, since the transient contribution
to \eqref{eq:window} is bounded in magnitude by
$2m|S_0-S_*|e^{-\beta t}/(\gamma\beta\Delta)$.

\section{What would select a common constant?}
For fixed bath mass, variance and clock, the plateau grows as $m+M$.
Achieving the same $H_*$ for different tracer masses requires
\begin{equation}\label{eq:universality}
 \nu_m=\frac{(m+M)s_m^2}{H_*}.
\end{equation}
This equation is a diagnostic of a proposed encounter law. Binary energy and
momentum conservation alone place no restriction of this form on the external
clock or reservoir preparation.

More explicitly, take a fixed admissible bath and a tracer initially at rest.
For any $0<\epsilon\le1$, replace every incoming velocity by
$W_n^{(\epsilon)}=\epsilon W_n$, keeping masses and clock unchanged.
Every elastic collision, centering assumption and speed bound remains valid.
Both $H_*$ and $\mathsf a(t)$ are multiplied by $\epsilon^2$.
Consequently this entire class admits arbitrarily small positive plateaus even
though every nondegenerate member relaxes to its own positive constant.
This is a countermodel test for these specified stochastic collision premises.

The next model should replace the supplied clock by a spatial encounter law
and test how it scales with reservoir velocity and density. The physical-time
relaxation mechanism is already explicit: persistent injection of velocity
variance balances collisional contraction. A universal quantum action scale
requires an additional physical relation that fixes this balance across systems.

Ben-Naim and Krapivsky provide a useful clock comparison: exponential decay
in collision number becomes algebraic decay in physical time in their cooling
model \cite[p.~6, (14)--(17)]{BenNaimKrapivsky2003InelasticMaxwell}.

\section{A spatial receiver with an exact collision clock}\label{sec:spatial}
An infinite equal-mass hard-point gas realizes the two-speed action plateau
with every particle's speed bounded by $u<c$. Its encounter clock follows
from the initial spatial statistics. This is the dichotomic Jepsen-gas case
analyzed by Balakrishnan, Bena and Van den Broeck
\cite[pp.~7--8, (11)--(12), pp.~17--19]{BalakrishnanBenaVanDenBroeck2002}.
The following construction makes its clock explicit.

Fix $m,\rho,u>0$. Put a tagged particle at the origin and independently put
Poisson particles of intensity $\rho$ on the rest of the line. Give every
particle an independent sign, with velocities $+u$ and $-u$ equally likely.
This is the Palm preparation: a typical particle is placed at zero, with
independent gas on either side. Particles move freely between instantaneous
elastic collisions, exchanging velocities and preserving their spatial order.
Each collision conserves momentum and kinetic energy. The initial energy
per unit length has mean $\rho m u^2/2$; the reservoir has infinite extent.

The equivalent free-trajectory construction lets particles cross and exchanges
their physical labels at crossings. Each marked Poisson stream is translated
rigidly, so the untagged homogeneous marked-gas law is invariant. A tagged
path on $[0,T]$ lies inside $[-uT,uT]$; only free lines starting in
$[-2uT,2uT]$ can intersect that region during this interval. There are almost
surely finitely many such lines and crossings. This constructs the local
dynamics without an accumulation of collisions. Positions are continuous
and piecewise linear; velocities are right-continuous with impulsive changes.

\begin{proposition}[Spatially generated reversals]\label{prop:spatial}
In this preparation the tagged velocity is a stationary symmetric two-state
Markov chain with reversal rate $\lambda=\rho u$. Consequently
\begin{align}
 C(s)&=u^2e^{-2\rho u s},\\
 \mathsf h(\Delta)&=\frac{mu}{\rho}
 \left[1-\frac{1-e^{-2\rho u\Delta}}{2\rho u\Delta}\right],
 & H_*&=\frac{mu}{\rho}>0.
 \label{eq:spatial-plateau}
\end{align}
Here $\mathsf h(\Delta)=m\Var[X(t+\Delta)-X(t)]/\Delta$,
and the limit is $\Delta\to\infty$ at fixed $m,\rho,u$ in the infinite gas.
\end{proposition}
\begin{proof}
Condition first on tagged initial velocity $+u$. Write $B_1<B_2<\cdots$
for the initial positions of negative-velocity free lines to its right, and
$-A_1>-A_2>\cdots$ for positive-velocity lines to its left; set $A_0=B_0=0$.
Independent thinning makes $(A_j)$ and $(B_j)$ independent Poisson arrival
sequences of intensity $\rho/2$. Their successive gaps are independent
exponentials of that rate.

The tag first switches from the line $ut$ to $B_1-ut$, then to $-A_1+ut$,
then to $B_2-ut$, and so on. Lines with the same velocity remain parallel;
at each switch the next opposing line is the next unused member of the
corresponding ordered sequence. Solving these intersections gives
\[
 t_{2j-1}=\frac{B_j+A_{j-1}}{2u},\qquad
 t_{2j}=\frac{B_j+A_j}{2u}.
\]
Thus successive waiting times alternate between
$(B_j-B_{j-1})/(2u)$ and $(A_j-A_{j-1})/(2u)$: all are independent
exponentials of rate $\rho u$. Reflection proves the same law for initial
velocity $-u$. The independent uniform initial sign is the stationary law
of this chain. Its covariance follows by taking the parity of its Poisson
reversal count. Integrating
$2m\int_0^\Delta(1-s/\Delta)C(s)\,ds$ proves the action formula.
\end{proof}

The mechanical statement agrees with the source's covariance (12). Its
discussion distinguishes a tag initially at $\pm u$ from a tag initially at
rest: p.~19 gives the ballistic term $e^{-\rho ut}\delta(X\mp ut)$ in the
former case, and a Bessel-weighted atom at zero in the latter. Our proposition
uses the former preparation throughout.

\subsection{The scale supplied by spatial preparation}
The mean collision duration is $1/(\rho u)$, and the plateau has the
action units of momentum times mean spacing: $H_*=mu/\rho$.
At fixed density, cooling $u\mapsto\epsilon u$ changes the collision rate
by $\epsilon$ and the plateau by $\epsilon$, rather than the quadratic
scaling of the independently clocked bath. At fixed $u$, increasing
$\rho$ also closes the plateau. Each family retains the hard speed ceiling.
A shared value $K$ across these gases therefore requires
$\rho_m=mu_m/K$; a preparation law would have to select this relation.

This realization advances the mechanical-clock obligation and inherits the
two-speed bridge and cut results through equality of the tagged path laws.
The positive plateau concerns long observation windows; microscopic refinement
still gives $\mathsf h(\Delta)\sim mu^2\Delta\to0$. The next test holds
$m,u,\rho$ fixed while changing spacing correlations, to identify how much of
the plateau is supplied by Poisson preparation. This is a classical impulsive
gas in a chosen frame, with a kinematic speed ceiling; a relativistic collision
theory or a smooth-force realization would add further dynamics.

\section{Return to the partition question}
The project's central question concerns a sequence of time partitions
$\pi_N=\{0=t_0^{(N)}<\cdots<t_N^{(N)}=T\}$ with mesh tending to zero.
The present relaxation law describes preparation time under a supplied bath.
It therefore serves as a comparison mechanism for the partition problem:
one must derive, from admissible cut-point refinements, both the quantity
that survives and the condition selecting its positive value. The collision
calculation isolates how a positive constant can instead enter through
reservoir energy and a clock. The next partition test tracks these inputs
explicitly when a refinement law is proposed.

\bibliographystyle{plain}
\bibliography{references/library}
\end{document}
